\Subhead{2.9.2}{Взаимодействия плоских волн в диспергирующих средах. Вторая гармоника. Параметрическое усиление и генерация}

Квадратичная поляризация $P^{(2)}\propto E^2$ смешивает частоты. Для $\omega_3=\omega_1+\omega_2$ эффективное накопление требует $\Delta k=k_3-k_1-k_2\simeq0$. Усреднение по быстрым фазам оставляет только резонансные члены. Выберем амплитуды $a_j$ так, что $|a_j|^2$ пропорционально потоку квантов; при точном синхронизме и вещественной связи
\begin{equation}\label{eq:threewave}
 a'_1=\ii\kappa a_2^*a_3,\quad
 a'_2=\ii\kappa a_1^*a_3,\quad
 a'_3=\ii\kappa a_1a_2.
\end{equation}
Умножение на сопряжённые амплитуды даёт $\dd|a_1|^2/\dd z=\dd|a_2|^2/\dd z=-\dd|a_3|^2/\dd z$. Это соотношения Мэнли–Роу; вместе с частотным условием они обеспечивают сохранение энергии.

При генерации второй гармоники $\omega_2=2\omega_1$ в приближении заданного поля $a_1$ имеем $a_2'\propto\ii a_1^2\e^{-\ii\Delta kz}$. Интегрирование даёт
\begin{equation}
 I_2(L)\propto I_1^2 L^2\operatorname{sinc}^2(\Delta kL/2).
\end{equation}
При несовпадении фаз вклады разных участков компенсируются; длина когерентности $L_c=\pi/|\Delta k|$. Квадратичный рост прекращается при истощении основного поля, потерях или уходе импульсов друг относительно друга.

Если высокочастотная волна $a_3=A_p$ задана, из \eqref{eq:threewave} следует $a_1''=|\kappa A_p|^2a_1$: сигнальная и холостая волны растут гиперболически. При расстройке пространственный инкремент $g=\sqrt{|\kappa A_p|^2-(\Delta k/2)^2}$. Для генератора нужна обратная связь и превышение потерь; усилитель с входным сигналом не требует резонаторного порога самовозбуждения.
